% Complete proof after independent actual-file audits.
\subsubsection*{Statement B.2u (Proved): weak upper-cone stabilization over every field}
\begin{theorem}
\label{B:B2-arbitrary-field-stabilization}
Let $k$ have characteristic $p>0$, let $K/k$ be finitely generated,
and suppose its class of proper geometrically normal models is
nonempty.  Given any such model $X_*$, there is a proper
geometrically normal modification $X_0\to X_*$ such that
\[
 \operatorname{Pic}^{\tau}_{X_0/k}\xrightarrow{\ \sim\ }
 \operatorname{Pic}^{\tau}_{W/k}
\]
for every proper geometrically normal modification $W\to X_0$.
This is the weak upper-cone assertion.  A constant domination-cofinal
diagram can still fail, as B.2q proves.
\end{theorem}

\begin{proof}
We first assume $K/k$ regular.  Choose the equivariant alteration
data of Lemma~\ref{B:equivariant-constant-data} over $X_*$:
a smooth projective geometrically integral $Y/l$, with $l/k$
finite, function field $L$, and a finite semilinear group
$\Gamma$ such that $L^\Gamma/K$ is finite purely inseparable.
Put
\[
 d=[L:K],\qquad P=\operatorname{Pic}^{\tau}_{Y/l},\qquad
 R=\operatorname{Res}_{l/k}P,\qquad A=R^\Gamma.
\]
For each permitted $X\to X_*$, normalize the main graph
$Y_X$ between $Y$ and $X$.  It has function field $L$, and
$Y_X\to Y$ is proper birational with normal source.  By B.1g
over $l$ its torsion Picard group is canonically $P$.
This identification respects the semilinear $\Gamma$-action.
Consequently B.2m gives compatible homomorphisms
\[
 h_X:P_X:=\operatorname{Pic}^{\tau}_{X/k}\longrightarrow A
           \subset R
\]
whose kernels are finite and killed by $d$.  The assertion
about finite kernels includes all parameter schemes, by the
fixed-Jacobian bound in B.2m.  Pullbacks among the $P_X$ are
closed immersions.  The group $R$ is of finite type and has
no geometric torus by B.2r(iii), which allows arbitrary finite
constant extensions.

Choose $X_1\to X_*$ maximizing $\dim P_X$.
The dimensions are bounded by $\dim R$, because the kernels
of $h_X$ are finite.  Let $A_*\subset P_{X_1}$ be its largest
abelian subvariety, furnished by B.2r.  Properness of $P_{X_1}$
implies that $P_{X_1}/A_*$ is finite.  Every subsequent
$W\to X_1$ has this same abelian subvariety under pullback,
since its dimension is maximal; we identify those copies of
$A_*$ without identifying the entire $P_W$.
Let
\[
                         A_0=h_{X_1}(A_*)\subset A.
\]
This is an abelian variety isogenous to $A_*$, since the
kernel is finite.  The distinction between $A_*$ and its
image $A_0$ is needed below.

\paragraph{Removing excessive abelian directions in the ambient group.}
Initially set $H=A$.  We will impose finitely many closed
subgroup conditions on $H$ and replace the current permitted
model by a domination, always keeping every subsequent image
$h_W(P_W)$ inside $H$.  Let $B\subset H$ be its largest
abelian subvariety, which exists by B.2r and contains $A_0$.
If $B=A_0$, no further condition is needed.
Otherwise $\dim B>\dim A_0$.

Lemma~\ref{B:invariant-abelian-realization} supplies an
abelian torsor $Q/k$ and a rational map $X_1\dashrightarrow Q$
whose pullback on connected Picard groups, evaluated in $R$,
has scheme-theoretic image exactly $B$.
Let $G$ be the normalization of the closed graph.
If it were geometrically normal, it would be a permitted
modification and its image $h_G(P_G)$ would contain $B$.
This contradicts maximality of $\dim P_{X_1}=\dim A_0$.
Thus $G$ is not geometrically normal.

Exactly as in B.2s, normality gives geometric $S_2$, and
regularity of $K/k$ gives geometric integrality.  Hence there
is a geometrically bad divisorial valuation $v$ on $G$.
Its centre on every permitted model has dimension at most
$n-2$, where $n=\operatorname{trdeg}_k K$, because a proper
birational normal morphism is an isomorphism near a
codimension-one point.  The divisor on $G$ has image of
dimension $n-1$ in $X_1\times Q$.
Choose a permitted $X_2\to X_1$ on which the centre $C_2$
of $v$ has maximal dimension.  For every later $W\to X_2$,
its centre $C_W$ has that same dimension, and
\[
                     F=k(C_2)\ \subset\ k(C_W)
\]
is a finite extension inside $\kappa(v)$.

Choose an extension $w$ of $v$ to the finite field extension
$L/K$.  It is a divisorial valuation of $L$.
Normalize $G$ in $L$ and take a normal proper common graph
$Z$ over this normalization, $Y$ and $X_2$, with function
field $L$.  By Chow's lemma followed by normalization, we
may and do take $Z$ projective over $k$
\cite[Tag 0200]{StacksB}; no geometric normality of $Z$ is
required.  The map $Z\to Y$ is birational, and the divisor
of $w$ persists because the proper birational map to the
normalization of $G$ is an isomorphism at its generic point.
Let $D$ be the normalization of that divisor on $Z$.
Then $D_F$ is proper normal integral over $F$, with function
field $\kappa(w)$, finite over $\kappa(v)$.
It is projective, so the relative Picard group and its closed
scheme-theoretic kernels have the usual projective
representability; the same conclusion also follows from
proper Picard representability over a field.

The restriction condition must be defined on all of $R$,
rather than merely on the known Picard images.  B.1g over
$l$ identifies $P$ with $\operatorname{Pic}^{\tau}_{Z/l}$.
The function field $L/l$ is regular, so $Z/l$ has a nonempty
smooth open.  Lemma~\ref{B:Pic-finite-constants-sheaf} gives
the full fppf identification
\[
 \operatorname{Pic}_{Z/k}
       \simeq\operatorname{Res}_{l/k}\operatorname{Pic}_{Z/l}.
\]
It follows that restriction gives a homomorphism
\[
                       R_F\longrightarrow\operatorname{Pic}_{D_F/F}.
\]
This uses no equality between torsion components under
restriction of scalars.  Let $K_w\subset R$ be the closed
$k$-subgroup defined by the vanishing of this restriction
after base change to $F$; its coefficient-ideal construction
and universal factorization property are exactly those
proved in B.2s.

All later images $h_W(P_W)$ are contained in $K_w$.
To check this, form a further normal common graph over $Z$
and $W$ and normalize the strict transform $D'$ of $D$.
The map $D'_F\to D_F$ is proper birational between normal
integral schemes, so its direct image of the structure sheaf
is $\mathcal O_{D_F}$.  Every parameter base change over the
field $F$ is flat, making relative Picard pullback injective
on all test schemes.  The map from $D'_F$ to $W$ factors
through $\operatorname{Spec}k(C_W)$.
Thus a line bundle from $W$ restricts to one pulled back
from the finite $F$-algebra $k(C_W)$.
Its relative Picard class is zero after the smooth
surjective generator-open cover of the parameter scheme.
This includes nonreduced parameters and inseparable centre
field extensions.  Injectivity on $D'_F\to D_F$ proves the
required vanishing on the fixed fibre.

The restriction does not vanish on all of $B$.
The image of $D$ in $X_1\times Q$ has dimension $n-1$:
the extension $\kappa(w)/\kappa(v)$ is finite and does not
alter the dimension of the original divisor image.
The image in $C_2\times Q$ dominates that image, while
$\dim C_2\le n-2$.
Therefore $D_F\to Q_F$ has positive-dimensional image.
After algebraic closure and normalization of a reduced
irreducible component, Albanese duality shows that its
pullback on $\operatorname{Pic}^0_Q$ is nonzero.
If $B$ were contained in $K_w$, this pullback would be zero,
because its map to $R$ has image $B$.  This contradiction
proves $B\not\subset K_w$.

Replace $H$ by $H\cap K_w$ and relabel the current model
$X_2$ as $X_1$.  The largest abelian subvariety of the new
$H$ is contained in the proper subgroup $B\cap K_w$, so
its dimension strictly drops, while it still contains $A_0$.
After finitely many steps we obtain a permitted model $X_3$
and a closed subgroup $H\subset A$ containing every later
$h_W(P_W)$ and having largest abelian subvariety exactly $A_0$.

\paragraph{A degree factor controls the remaining comparison kernel.}
By B.2r there is a fixed integer $N>0$ with
$[N]H\subset A_0$.
For every later $W$, we claim that
\[
                         [dN]P_W\subset A_*.
\]
Indeed, for a local section $x$ of $P_W$ the section
$N h_W(x)$ lies in $A_0$.
The isogeny $A_*\to A_0$ is fppf-surjective, so fppf locally
we can choose $a\in A_*$ with $h_W(a)=N h_W(x)$.
Then $Nx-a\in\ker h_W$, which is killed by $d$.
Consequently $dNx=da\in A_*$, proving the claim as a
group-scheme assertion.

The quotient $P_W/A_*$ is finite, since $A_*$ has the
maximal dimension of $P_W$.  It is killed by $dN$, and
\[
                  P_W[dN]\longrightarrow P_W/A_*
\]
is fppf-surjective: subtract a $(dN)$-division of
$(dN)x$ in the abelian variety $A_*$ from a local lift $x$.
After one fixed algebraically closed extension, B.2d embeds
all these $P_W$ into one $g$-dimensional Jacobian $J$;
normal graphs are required only over that algebraically
closed field.  Hence
\[
                    |P_W/A_*|\le |J[dN]|=(dN)^{2g}.
\]
Choose $X_0\to X_3$ with maximal finite order
$|P_{X_0}/A_*|$.  Every later closed immersion of Picard
groups induces a closed immersion between finite quotients
of that same order and is therefore an isomorphism.
This proves the theorem for regular $K/k$.

In general, the existence of a geometrically normal model
implies that the algebraic constant field $k_0$ of $K/k$
is finite separable over $k$, and that $K/k_0$ is regular.
We apply the theorem over $k_0$ to the model $X_*$.
Its underlying proper modifications are the same, and the
Picard isomorphisms over $k$ are obtained by
$\operatorname{Res}_{k_0/k}$, or checked after a finite
separable extension splitting the constants.  This completes
the arbitrary-field proof.
\end{proof}

\begin{remark}[The three formulations are distinct]
We have now proved both uniform $h^1$ boundedness and the
weak upper-cone assertion over arbitrary fields.  The
constant domination-cofinal formulation is false over an
imperfect field by B.2q, and the indicated model category
need not be filtered by B.2i.  The theorem does not replace
these counterexamples with a claim of cofinality: two
different stable upper cones may have nonisomorphic values.
\end{remark}
